% ============================================================
\section{Ejercicio 1}
% ============================================================

\begin{tcolorbox}[enunciado]
Se sabe que un algoritmo de ordenamiento Q con complejidad $O(f(n))$ y tiempo de procesamiento $T(n)=cf(n)$, tarda 0.4 segundos en ejecutarse cuando la lista a ordenar tiene 10,000 elementos.
\begin{enumerate}
    \item Estima el tiempo de ejecución para una lista de 60,000 elementos.
    \item Estima el tiempo de ejecución para una lista de 1,500,000 elementos.
    \item Determina el tamaño máximo de la lista que puede ser ordenada en menos de 10 segundos.
\end{enumerate}
Puedes asumir que el tiempo de ejecución del algoritmo Q tiene una complejidad $O(n \log n)$
\end{tcolorbox}

% ------------------------- Solución -------------------------
\subsection{Solución}

\begin{enumerate}
    %============================PUNTO 1============================
    \item 
    El \textbf{primer punto} pide estimar el tiempo de ejecución del algoritmo $Q$ cuando la lista a ordenar tiene $n = 60,000$ elementos.

    La complejidad del algoritmo es $O(n \log n)$ y el tiempo real se modela como:
    $$T(n) = c \cdot n \log(n)$$
    Utilizamos el caso base proporcionado para determinar el valor de $c$: con $n = 10,000$, el tiempo es $T(10,000) = 0.4$ segundos. 

    \subsubsection*{Elección de la Base del Logaritmo}

    Antes de continuar creemos necesario señalar que \textbf{la base que elijamos no afecta en absoluto a las predicciones de tiempo}. Esto ocurre porque el cambio de base de un logaritmo es una conversión lineal que la constante $c$ absorbe por completo.
    
    \textbf{Demostración formal:}\\
    Si cambiamos el logaritmo de una base $a$ a una base $b$, usamos la propiedad:
    $$\log_a(n) = \frac{\log_b(n)}{\log_b(a)}$$
    
    Si planteamos la ecuación con la base $a$:
    $$T(n) = c_a \cdot n \log_a(n) = c_a \cdot n \frac{\log_b(n)}{\log_b(a)} = \left( \frac{c_a}{\log_b(a)} \right) \cdot n \log_b(n)$$
    
    Definiendo una nueva constante $c_b = \frac{c_a}{\log_b(a)}$, la ecuación se convierte en:
    $$T(n) = c_b \cdot n \log_b(n)$$
    
    Al despejar la constante utilizando el caso base $T(10,000) = 0.4$ s, el valor de $c$ se ajustará automáticamente según la base que elijamos, y cualquier cálculo posterior para otro valor de $n$ dará exactamente el mismo tiempo en segundos.

    \subsubsection*{Encontrar el valor de la contante $c$}

    Usando logaritmo en base 10 (por simplicidad de cálculo manual).

    Sustituimos los valores en el modelo:
    $$0.4 = c \cdot 10,000 \cdot \log_{10}(10,000)$$
    $$0.4 = c \cdot (10,000 \cdot 4)$$
    $$0.4 = c \cdot 40,000$$
    $$c = \frac{0.4}{40,000} = 10^{-5} = 0.00001$$
    
    \subsubsection*{Estimar el tiempo para $n = 60,000$ elementos}
    Ahora que conocemos la constante $c$, evaluamos el tiempo de ejecución en $n = 60,000$.
    
    Sustituimos en la ecuación:
    $$T(60,000) = 10^{-5} \cdot 60,000 \cdot \log_{10}(60,000)$$
    $$T(60,000) = 0.6 \cdot \log_{10}(60,000)$$
    Calculamos $\log_{10}(60,000)$:
    $$\log_{10}(60,000) = \log_{10}(6 \cdot 10^4) = \log_{10}(6) + \log_{10}(10^4) \approx 0.77815 + 4 = 4.77815$$
    Multiplicamos:
    $$T(60,000) \approx 0.6 \cdot 4.77815 \approx \mathbf{2.867} \text{ \textbf{segundos}}$$
    
    
    \subsubsection*{Conclusión:}
    El tiempo estimado de ejecución para una lista de $60,000$ elementos es de aproximadamente \textbf{2.867 segundos} (o $2.87$ segundos si redondeamos a dos decimales).

    \vspace{0.4cm}
    %============================PUNTO 2============================
    
    \item 
    Para resolver el \textbf{segundo punto} debemos estimar el tiempo de ejecución del algoritmo $Q$ cuando la lista tiene: \textbf{$n = 1,500,000$ elementos}.
    
    Utilizaremos el valor de la constante $c$ que calculamos y la base logarítmicas 10.
    
    \subsubsection*{Valor de la constante $c$}
    
    Sabemos que el tiempo de ejecución del algoritmo se modela como:
    $$T(n) = c \cdot n \log(n)$$
    
    Donde, a partir de nuestro caso base ($T(10,000) = 0.4$ s) y utilizando logaritmo en base 10:
    $$c = 10^{-5} = 0.00001$$
    
    \subsubsection*{Estimar el tiempo para $n = 1,500,000$ elementos}
    Sustituimos $n = 1,500,000$ y $c = 10^{-5}$ en nuestra ecuación:
    $$T(1,500,000) = 10^{-5} \cdot 1,500,000 \cdot \log_{10}(1,500,000)$$
    
    Simplificamos el producto de la constante $c$ por $n$:
    $$10^{-5} \cdot 1,500,000 = 15$$
    
    Ahora, calculamos $\log_{10}(1,500,000)$. 
    $$\log_{10}(1,500,000) = \log_{10}(1.5 \cdot 10^6) = \log_{10}(1.5) + \log_{10}(10^6)$$
    Sabiendo que $\log_{10}(1.5) \approx 0.17609$ y $\log_{10}(10^6) = 6$:
    $$\log_{10}(1,500,000) \approx 0.17609 + 6 = 6.17609$$
    
    Multiplicamos por el factor restante:
    $$T(1,500,000) \approx 15 \cdot 6.17609 \approx \mathbf{92.641} \text{ \textbf{segundos}}$$
    
    \subsubsection*{Conclusión:}
    El tiempo estimado de ejecución para una lista de $1,500,000$ elementos es de aproximadamente \textbf{92.64 segundos} (o \textbf{1 minuto y 32.64 segundos}).
    
    Este resultado refleja el comportamiento de la complejidad $O(n \log n)$: aunque el tamaño de la lista aumentó \textbf{150 veces} (de $10,000$ a $1,500,000$), el tiempo de procesamiento solo aumentó unas \textbf{231 veces} (de $0.4$ s a $92.64$ s).

    \vspace{0.4cm}

    %============================PUNTO 3============================
    \item 
    Para resolver el \textbf{tercer punto} de manera exacta, utilizaremos la \textbf{función $W$ de Lambert \cite{herrera2025lambert}}. Este método evita las aproximaciones por tanteo y proporciona una solución analítica cerrada para la ecuación.
    
    \subsubsection*{1. Formulación de la Desigualdad}
    El objetivo es encontrar el tamaño máximo de la lista $n$ tal que el tiempo de ejecución del algoritmo sea estrictamente menor a 10 segundos:
    $$T(n) < 10$$
    
    A partir de los resultados obtenidos con el caso base $T(10,000) = 0.4$ s, podemos plantear la ecuación utilizando cualquier base logarítmica. 
    
    \textbf{Usamos Base 10 (donde $c = 10^{-5}$):}
    $$10^{-5} \cdot n \log_{10}(n) < 10 \implies n \log_{10}(n) < 10^6$$
    Cambiando el logaritmo a base natural ($\log_{10}(n) = \frac{\ln(n)}{\ln(10)}$):
    $$n \frac{\ln(n)}{\ln(10)} < 10^6 \implies n \ln(n) < 10^6 \ln(10)$$
    
    
    La ecuación que define el límite exacto es:
    $$n \ln(n) = 10^6 \ln(10)$$
    
    \subsubsection*{2. Transformación a la Forma de Lambert}
    La \textbf{función $W$ de Lambert} es la función inversa de $f(w) = w e^w$. Es decir:
    $$w e^w = z \iff w = W(z)$$

    Para llevar nuestra ecuación $n \ln(n) = 10^6 \ln(10)$ a esta forma exacta, aplicamos la identidad algebraica de la función exponencial, escribiendo $n$ como $e^{\ln(n)}$:
    $$e^{\ln(n)} \ln(n) = 10^6 \ln(10)$$
    
    Reordenando los términos del lado izquierdo:
    $$\ln(n) \cdot e^{\ln(n)} = 10^6 \ln(10)$$
    
    Esta ecuación ya tiene exactamente la estructura $w e^w = z$, donde:
    \begin{itemize}[leftmargin=*]
        \item La variable incógnita es $w = \ln(n)$
        \item El término constante es $z = 10^6 \ln(10)$
    \end{itemize}
    
    \subsubsection*{3. Aplicación de la Función $W$ de Lambert}
    Aplicamos la función $W$ (específicamente la rama principal $W_0$, ya que nuestro argumento $z$ es positivo) en ambos lados de la ecuación para despejar $\ln(n)$, es decir:

    \vspace{0.4cm}
    \begin{tcolorbox}[
        colback=gray!10!white,
        colframe=gray!60!black,
        fonttitle=\bfseries,
        title=Nota,
        arc=4pt,
        boxrule=0.5pt,
        left=6pt, right=6pt, top=4pt, bottom=4pt
    ]
    La \textbf{función $W$ de Lambert} se define matemáticamente con un único propósito: ser la función inversa de la expresión $f(x) = x e^x$. Es decir, $W$ es la herramienta diseñada para ``deshacer'' la operación de multiplicar una variable por su propia exponencial.
    
    Por la propiedad de cualquier función inversa, sabemos que al aplicar la inversa a la función original, estas se cancelan dejando únicamente el argumento:
    $$f^{-1}(f(x)) = x$$
    
    Si aplicamos esto a nuestra definición:
    \begin{enumerate}[leftmargin=*]
        \item La función directa es: \textbf{$f(x) = x e^x$}
        \item La función inversa es: \textbf{$f^{-1}(z) = W(z)$}
        \item Al componerlas, obtenemos la identidad fundamental de Lambert: 
    \end{enumerate}
    $$W(x \cdot e^x) = x$$
    \end{tcolorbox}

    $$W(\ln(n) \cdot e^{\ln(n)}) = W(10^6 \ln(10))$$
    $$\ln(n) = W(10^6 \ln(10))$$
    
    Para despejar finalmente $n$, aplicamos la función exponencial (base $e$) a ambos miembros:
    $$n = e^{W(10^6 \ln(10))}$$
    
    \textbf{Simplificar la expresión}\\
    Existe una propiedad fundamental de la función $W$ de Lambert que establece que:
    $$W(z) \cdot e^{W(z)} = z \implies e^{W(z)} = \frac{z}{W(z)}$$
    
    Sustituyendo esta propiedad en nuestra solución para $n$ (donde $z = 10^6 \ln(10)$), reducimos la expresión exponencial a una división mucho más sencilla:
    $$n = \frac{10^6 \ln(10)}{W(10^6 \ln(10))}$$
    
    \subsubsection*{4. Evaluación Numérica de la Solución}
    Ahora evaluamos los valores numéricos:
    \begin{enumerate}[leftmargin=*]
        \item Calculamos el argumento $z$:
        $$z = 10^6 \ln(10) \approx 1,000,000 \cdot 2.30258509 = 2,302,585.09$$
        \item Evaluamos la función $W_0$ de Lambert para este argumento. Dado que es una función trascendente, su valor se obtiene mediante aproximaciones numéricas de alta precisión (por ejemplo, con métodos de Newton-Raphson, librerías científicas como SciPy o en este caso utilizamos WolframAlpha \cite{pensando2021lambert}):
        $$W_0(2,302,585.09) = ProductLogW(10^6 \ln(10)) $$
        $$W_0(2,302,585.09) \approx 12.151944$$
        \textbf{Nota:} La función $W$ de Lambert aparece como ProductLog() en \textbf{WolframAlpha}.
        \item Calculamos el valor final de $n$:
        $$n \approx \frac{2,302,585.09}{12.151944} \approx 189,481.28$$
    \end{enumerate}

    % \begin{figure}[H]
    %     \centering
    %     \includegraphics[width=0.8\linewidth]{ejercicio1_punto3_resultado_final.png}
    %     \caption{Captura de pantalla de la obtención del resultado final}
    % \end{figure}
    
    \subsubsection*{5. Conclusión}
    Dado que $n$ representa el número de elementos físicos de una lista, este debe ser necesariamente un \textbf{número entero}. 
    
    \begin{itemize}[leftmargin=*]
        \item Para $n = 189,481$, el tiempo estimado de ejecución es de \textbf{9.99998 segundos} (cumple estrictamente con $T(n) < 10 \implies n \log_{10}(n) < 10^6$).

        % \begin{figure}[H]
        %     \centering
        %     \includegraphics[width=0.6\linewidth]{1_03_primera_aproximacion.png}
        %     \caption{Captura de pantalla de la primera aproximación numérica}
        % \end{figure}

        \item Para $n = 189,482$, el tiempo estimado de ejecución es de \textbf{10.00004 segundos} (excede el límite de 10 segundos).

        % \begin{figure}[H]
        %     \centering
        %     \includegraphics[width=0.6\linewidth]{1_03_segunda_aproximacion.png}
        %     \caption{Captura de pantalla de la segunda aproximación numérica}
        % \end{figure}
        
    \end{itemize}
    
    Por lo tanto, la solución exacta y justificada analíticamente para el tamaño máximo de la lista es:
    $$\mathbf{n = 189,481 \text{ \textbf{elementos}}}$$

\end{enumerate}
